% !TEX root = userguide.tex

\def\headdate{12th November 2021}

\section{Mixed Stokes}
\label{sec:mixedstokes}

In this section the mixed Stokes formulation is introduced and some examples
are given.

\subsection{Mixed Stokes formulation.}

We consider the mixed Stokes formulation in a region $\Omega$:
Find $(\vek u,p,\ten\tau)$ such that,
\begin{alignat*}{2}
      -\nabla\cdot(2\eta_\text{s}\ten D)+\nabla p-\nabla\cdot\ten \tau
        &=\vek f\qquad&&\text{in }\Omega\\
\nabla\cdot\vek u&=0\qquad&&\text{in }\Omega\\
 \ten\tau-2\eta\ten D&=\ten 0\qquad&&\text{in }\Omega
\end{alignat*}
where $\ten D=(\nabla\vek u+\nabla\vek u^T)/2$, $\eta_\text{s}$ is the direct (solvent) viscosity
and $\eta$ is the mixed (polymer) viscosity. The boundary conditions are
\begin{gather*}
    \vec u = \vec u_D\text{ on }\Gamma_D\\
    -p\vec n+2\eta_\text{s} \ten D\cdot\vec n+\ten \tau\cdot\vek n=\vec t_{\text{N}}
\text{ on }\Gamma_{\text{N}}
\end{gather*}
where the boundary $\Gamma=\Gamma_D\cup\Gamma_{\text{N}}$ has been split into a
Dirichlet and a Neumann part. The Neumann condition is equivalent
to an imposed traction of $\vec t_{\text{N}}$.

The weak form becomes: Find $(\vek u,p,\ten\tau)$ such that,
\begin{alignat*}{2}
  \big((\nabla\vek v)^T,2\eta_s\ten D+\ten \tau\big)-(\nabla\cdot\vek v,p)&=(\vek v,\vek t_\text{N})_{\Gamma_\text{N}}+(\vek v,\vek f)
&&\qquad\text{for all }\vek v \\
  (q,\nabla\cdot\vek u)&=0,&&\qquad\text{for all }q\\
(\ten s,\ten\tau -2\eta\ten D)&=0,&&\qquad\text{for all }\ten s
\end{alignat*}
or
\begin{alignat*}{2}	
  \big(\ten D_v,2\eta_s\ten D\big)-(\nabla\cdot\vek v,p)
+ \big(\ten D_v,\ten \tau\big)&=(\vek v,\vek t)_\Gamma+(\vek v,\rho\vek b)
&&\qquad\text{for all }\vek v \\
  -(q,\nabla\cdot\vek u)\hspace{3.5cm}&=0,&&\qquad\text{for all }q\\
\big(\ten s,\ten D\big)\hspace{2cm}-\frac1{2\eta}(\ten s,\ten\tau) &=0,&&\qquad\text{for all }\ten s
\end{alignat*}
In this equation
$\ten D_v$ is defined by
\[
\ten D_v=(\nabla\vek v+\nabla\vek v^T)/2
\]
Notes:
\begin{itemize}
   \item System is symmetric. For that, we multiplied the continuity equation and the Newtonian constitutive equation for $\ten\tau$ with -1 and also divided the latter by $2\eta$.
   \item The weak form is equivalent to the saddle point problem: find the
         stationary value of the functional
\[
         J(\vek u,p,\ten\tau)=
    \frac122\eta_s\big(\ten D,\ten D)\big)-(\nabla\cdot\vek u,p)
   +\big(\ten D,\ten \tau\big)-\frac 12\frac1{2\eta}(\ten\tau,\ten\tau)
\]
\item
Compatibility of the three-field system $(\vek u_h,p_h,\ten\tau_h)$ requires \cite{baranger_finite_1992,fortin_convergence_1989}:
\begin{enumerate}
   \item The $(\vek u_h,p_h)$ spaces need to be compatible for standard Stokes
(inf-sup or LBB condition for $(\vek u_h,p_h)$).
  \item For $\eta_s=0$ (no solvent) an \emph{extra} inf-sup condition for the
$(\vek u_h,\ten\tau_h)$ combination must be fulfilled.
\end{enumerate}
  \item The first compatibility condition is easy: use a standard compatible $(\vek u_h,p_h)$
    element.
  \item For $\eta_s\neq0$ the extra inf-sup compatibility condition  for
    $(\vek u_h,\ten\tau_h)$ is \emph{not necessary} in theory. However,
    in practise when the mixed Stokes is used for viscoelastic fluid flow,
    $\eta_s$ can be quite small (for polymer melts $\eta_s=0$) and stabilization
    due to $\eta_s$ is not enough in that case and the extra inf-sup
    condition needs to be fulfilled.
  \item The condition for $(\vek u_h,\ten\tau_h)$ is not easy. The matrix system
    for $\eta_s=0$ looks like:
\begin{center}
\includegraphics{mat2}
\end{center}
The $(\vek u_h,\ten\tau_h)$ parts look like a mixed system with a zero block.
This means that at least 
\[
    N_\tau\geq N_u
\]
i.e.\ the stress space must be \emph{larger} than the velocity space (``the
number of gradients must be larger than the primary variable'').
\item Compatible mixed Stokes elements are, for example, the Marchal \& Crochet $4\times 4$ $Q_1$-subelement \cite{marchal_new_1987} and the $Q_3$ element suggested by Fortin \& Pierre \cite{fortin_convergence_1989}.
\item Compatibility can also be achieved using the DEVSS technique (see Sec.~\ref{sec:weakviscoelastic}), but requires an additional field variable.
\end{itemize}

\subsection{Stokes flow through an abrupt contraction using mixed Stokes
 elements.}
\label{sec:mixedcontraction}

The example files can be obtained by typing at the command line:
\begin{quote}
  \texttt{getexample mixed\_stokes}
\end{quote}

The first example \texttt{mixed\_stokes1.f90} uses a $Q_2/Q_1$ Taylor-Hood
interpolation for velocity/pressure and a $Q_p,\, p\geq 1$ (continuous) interpolation for the stress.

The second example \texttt{mixed\_stokes2.f90} uses a $Q_2/Q_1$ Taylor-Hood
interpolation for velocity/pressure and a $Q_1\text{-iso-}Q_p,\, p\geq 1$ (continuous) interpolation for the stress. For $p=4$, this is the mixed Stokes element of Marchal \& Crochet \cite{marchal_new_1987}.

